\section{From Penance to Viaticum}

The Church accompanies the dying through the sacramental life in which she has nourished them from the beginning. In Penance, Christ reconciles the sinner. In Anointing of the Sick, he strengthens the person who is seriously ill or weakened by old age. In the Eucharist as Viaticum, he gives himself as food for the final journey. Through these sacraments the dying receive his forgiveness, strength, and presence as they approach the end of earthly life (CCC 1524--1525).

Preparation therefore includes learning to receive forgiveness before our last illness. An honest confession requires us to speak without arranging an excuse at the center of every admission. We name the sin because we desire freedom from it. The habit of approaching the sacrament can gradually overcome the fear that being fully known means being rejected. Christ already knows the truth, and his mercy calls us to bring the truth willingly before him. Shame need not be allowed to reserve a part of our life from absolution.

The thought of death can uncover something long withheld. A person may have carried a grave wrong for years, hoping that silence would diminish its reality. Another may have become distant from the sacraments through a combination of shame, neglect, and the assumption that returning would be impossibly difficult. The present invitation is simple in its direction, even when the circumstances need patient help: seek a priest and begin to speak truthfully. There is no need to settle every anxiety about the conversation before requesting it. The desire to return is already something to bring.

The sacrament of Anointing need not be reserved for the last few moments. The fitting time has arrived when a member of the faithful begins to be in danger of death through sickness or old age. The Catechism also identifies serious operations and increasing frailty as occasions for seeking it, and permits its repetition when a new grave illness arises or the condition becomes more serious (CCC 1514--1516). The sick person and those caring for him should contact a priest early enough for the Church's ministry to be received with whatever attention and freedom the condition allows. Waiting for certainty about the final hour can deprive a person of that opportunity.

James directs the sick to call for the priests of the Church, who pray and anoint in the Lord's name (Jas 5:14--15). The Catechism describes the sacrament's gifts as strengthening, peace, and courage, renewed trust in God, and closer union with Christ's Passion (CCC 1520--1523). Bodily healing is included within its purpose if God wills it. Where illness continues, the sick person still receives grace to endure it in union with Christ.

Anointing also places suffering within the Church's communion. The sick person receives her intercession and, through union with Christ, contributes to her sanctification (CCC 1522). A life whose outward activities have diminished can still bear fruit for others. This truth should be offered gently, with respect for the sufferer's freedom and strength. We cannot demand that another person make visible spiritual use of every pain. The sacrament gives what the weakened person needs; it is not a test designed to prove how well he is coping.

Viaticum brings the Eucharist to the threshold of death. The Lord who once gave his life now gives his risen body as the pledge of resurrection. The promise that he will raise those who eat his flesh and drink his blood appears at John 6:55 in the Douay--Rheims, corresponding to John 6:54 as cited by CCC 1524. The dying Christian receives the same Lord who has nourished the Church throughout the pilgrimage. The one who can no longer come to the assembly remains a member of Christ's body, visited and fed by its Head.

In receiving Communion through life, we have been receiving the Lord whom we hope to receive at death. Each reception calls for faith in his presence, gratitude for his gift, and a willingness to live according to his charity. We can prepare for Viaticum by attending to those dispositions now. If illness later makes concentration difficult, the habits formed in ordinary worship may help us turn toward him with the attention still possible. We can also make our desire for Communion known while we are able to express it clearly.

Practical preparation serves this sacramental readiness. We can make our wish for a priest known to family and caregivers, keep parish contact information accessible, and speak plainly about the importance of confession, anointing, and Communion when serious illness comes. The particulars will depend on the person's condition and the priest's pastoral judgment. These arrangements are acts of love toward those who may otherwise have to guess our wishes while distressed. We are providing room for the Church to do what she has received from Christ, and for ourselves to receive it as fully as possible.

